% Capitulo: Dados como código
\chapter{Dados como código}
\label{cap:dados-codigo}

Durante muito tempo, o trabalho de dados viveu fora das práticas que a engenharia de
software levou décadas para construir. Consultas ficavam salvas no computador de quem as
escreveu, a lógica de transformação estava escondida dentro da ferramenta de BI, e a versão
\enquote{oficial} de um painel era a que estava publicada, sem histórico de como chegou até
ali. Quando algo dava errado, ninguém sabia o que tinha mudado. Esta parte do livro trata de
trazer para os dados a disciplina da engenharia de software, e começa pela base: tratar
transformações, testes e, sempre que possível, painéis como código.

O movimento tem nome e uma literatura própria. A cultura DevOps nasceu para aproximar
desenvolvimento e operação de software, e se apoia em três ideias: fazer o trabalho fluir
em pequenos lotes do desenvolvimento à produção, ter retorno rápido sobre o que dá errado, e
aprender continuamente com isso \parencite{kim2016}. As três se aplicam aos dados quase sem
adaptação.

\section{Por que versionar}
\label{sec:versionar}

Versionar é guardar cada mudança no código com a informação de quem mudou, quando e por
quê, num sistema que permite ver o histórico, comparar versões e voltar atrás. A ferramenta
padrão para isso é o Git, e há um livro completo e gratuito sobre ela \parencite{chacon2014}.
Para quem vem da análise, o mínimo necessário cabe em poucos comandos:

\begin{terminal}
git switch -c receita-por-cliente         # cria um ramo para a mudança
# edita o modelo e o teste correspondente
git add modelos/ testes/
git commit -m "Receita por cliente com frete no grão do pedido"
git push -u origin receita-por-cliente    # envia e abre o pedido de revisão
\end{terminal}

Esse pequeno ritual muda muita coisa. Quando um número muda de um dia para o outro, o
histórico mostra qual alteração causou a mudança. Quando uma alteração quebra algo, voltar à
versão anterior leva minutos. Quando duas pessoas mexem no mesmo modelo, o sistema mostra o
conflito em vez de uma sobrescrever o trabalho da outra. E quando alguém novo chega ao time,
o histórico conta como cada regra de negócio chegou ao formato atual.

O que deve ser versionado vai além das consultas: os modelos de transformação, os testes, as
definições de métricas, a configuração dos pipelines, a documentação e, quando a ferramenta
permite, os próprios painéis. Muitas ferramentas de BI já conseguem exportar a definição dos
painéis como arquivos de texto, o que permite tratá-los com as mesmas práticas. Onde isso
não é possível, vale ao menos manter a lógica fora do painel, na camada de transformação,
onde ela pode ser versionada e testada.

Com a IA, o versionamento ganha uma importância nova. Quando parte do código é escrita por
um assistente ou por um agente, o histórico é o registro de cada alteração e de quem a
aprovou. É ele que permite responder, semanas depois, por que aquele filtro está ali.

\section{Revisão de código}
\label{sec:revisao-codigo}

Revisão de código é a prática de uma segunda pessoa ler cada alteração antes que ela chegue à
produção. Na forma moderna, acontece pelo pedido de revisão (\emph{pull request}): a mudança
fica num ramo separado, o revisor lê as diferenças, comenta, pede ajustes e só então aprova a
integração ao código principal.

O motivo mais citado para revisar código é encontrar defeitos. Um estudo com desenvolvedores
de uma grande empresa de software mostrou que, embora encontrar defeitos fosse a principal
expectativa, boa parte do valor real da revisão estava em outras coisas: transferência de
conhecimento, consciência do time sobre o que está mudando, e soluções alternativas que o
autor não tinha considerado \parencite{bacchelli2013}. Em times de dados, isso é especialmente
verdadeiro. A revisão é muitas vezes o momento em que alguém que conhece o negócio percebe que
a regra implementada não é a regra que a área usa.

Revisar SQL tem algumas perguntas próprias, que valem mais do que qualquer guia de estilo:

\begin{itemize}
  \item Qual é o grão de cada tabela que entra nesta consulta, e qual é o grão do resultado?
  \item Alguma junção pode multiplicar linhas (\veja{sec:consulta-errada})?
  \item Os filtros fazem sentido para o negócio? Algum registro legítimo está ficando de fora?
  \item Os nomes de tabelas e colunas são claros para quem não escreveu o código?
  \item Há testes que garantem as regras importantes, e eles passam?
\end{itemize}

Com a IA, a revisão vira o principal ponto de controle do processo. Quando o código é gerado
em segundos, a etapa que limita a qualidade deixa de ser a escrita e passa a ser a leitura.
Isso tem duas consequências práticas. A primeira é que o autor de uma alteração gerada por IA
continua sendo o responsável por ela, e deve revisá-la antes de pedir a revisão de outra
pessoa. A segunda é que alterações grandes, geradas de uma vez, são difíceis de revisar bem.
Vale pedir à IA alterações pequenas e focadas, do mesmo jeito que se pediria a um colega.

\antesdepois{A revisão era uma etapa entre muitas. O autor escrevia o código devagar,
entendia cada linha, e o revisor conferia principalmente estilo e regras de negócio.}
{A revisão é o principal controle de qualidade. O código chega rápido e em volume, e o revisor
precisa conferir também se o autor entendeu o que a IA escreveu.}

\section{Desenvolvimento e produção}
\label{sec:ambientes}

Testar em produção é a forma mais rápida de descobrir que algo estava errado, e a mais cara.
Um modelo alterado direto no ambiente que alimenta os painéis da diretoria pode publicar
números errados por horas antes de alguém perceber. A solução, conhecida há décadas na
engenharia de software, é separar ambientes.

O mínimo é ter dois: desenvolvimento, onde cada pessoa trabalha e testa as próprias
alterações, e produção, onde ficam os dados que as pessoas usam para decidir. Em armazéns de
dados na nuvem, isso costuma se traduzir em esquemas ou bancos separados, com a mesma
estrutura e permissões diferentes. Cada desenvolvedor pode ter o seu espaço de trabalho, com
uma cópia ou uma amostra dos dados de produção.

Um terceiro ambiente, de integração, completa o desenho. Nele, cada pedido de revisão é
construído e testado de forma isolada, antes de qualquer integração ao código principal. É o
que permite que o revisor veja não só o código, mas também o resultado: as tabelas geradas,
os testes que passaram e os que falharam. O capítulo~\ref{cap:ci-cd} trata de como
automatizar esse caminho.

Dois cuidados acompanham a separação de ambientes. O primeiro são as credenciais: senhas e
chaves de acesso nunca entram no código versionado, e cada ambiente usa as suas, com o
mínimo de permissão necessário. O segundo são os dados sensíveis: ambientes de
desenvolvimento costumam ter controles mais frouxos, e dados pessoais copiados de produção
para lá precisam ser mascarados (\veja{sec:seguranca}).

\rascunho{
  Relato: como foi introduzir versionamento, revisão e ambientes separados num time que
  não tinha essas práticas (resistências, o que convenceu as pessoas, o que mudou nos
  resultados).
}

\begin{leituras}
  \leitura{kim2016}{O manual da cultura DevOps. Os princípios de fluxo, retorno rápido e
    aprendizado contínuo se aplicam diretamente a times de dados.}
  \leitura{chacon2014}{O livro de referência sobre Git, completo e disponível gratuitamente
    na internet.}
  \leitura{bacchelli2013}{O estudo que mostra que o valor da revisão de código vai muito além
    de encontrar defeitos.}
\end{leituras}
